% GENERATED FILE. Edit canonical lessons, then run npm run generate:books.
% canonical-source-hash: fnv1a64:45471ed2efe50409
% canonical-lessons: JA-C261-ni-yorimasu, JA-C261-baai, JA-C261-kau-deshou, JA-W261-six-lines, JA-R261-it-depends

\chapter{It Depends, and What Would Happen}
\label{ch:ja-it-depends-and-what-would-happen}
\hlchaptermodality{261}

\begin{chapteropening}
\textbf{By the end of this chapter:} \emph{I can say what something depends on (tenki ni yorimasu), say what is done in a given case (ame no baai wa), say what would happen if (okane ga attara, kau deshō), and tell a trip in six connected lines.}

Complete JA-R261-it-depends: say what a plan depends on and what would happen if.
\end{chapteropening}

\section[\texorpdfstring{\ja{〜}\ja{に} \ja{よります} (ni yorimasu)}{〜に よります (ni yorimasu)}]{\ja{〜}\ja{に} \ja{よります} (ni yorimasu) — it depends on}

\label{lesson:JA-C261-ni-yorimasu}



\begin{quote}
Say \textquotedblleft{}if it is fine tomorrow, I will go to the mountains\textquotedblright{}. (\emph{moshi ashita haretara, yama ni ikimasu}.) Then \textquotedblleft{}the weather\textquotedblright{}. (\emph{tenki}.)
\end{quote}

\subsection*{You'll want to know: \ja{〜}\ja{に} \ja{よります}}
\begin{quote}
\textbf{\ja{てんきに} \ja{よります。}}
\end{quote}

\begin{quote}
\emph{tenki ni yorimasu.} — \textquotedblleft{}It depends on the weather.\textquotedblright{}
\end{quote}

What it depends on takes \textbf{\ja{に}}, and \textbf{\ja{よります}} (\emph{yorimasu}) follows. This \textbf{\ja{よる}} is \textquotedblleft{}to depend on\textquotedblright{}, a different verb from the \textbf{\ja{よる}} of \textbf{\ja{よります}} \textquotedblleft{}to drop in\textquotedblright{}, though the two sound the same; both are godan verbs.

\begin{quote}
\textbf{\ja{それは} \ja{じかんに} \ja{よります。}}
\end{quote}

\begin{quote}
\emph{sore wa jikan ni yorimasu.} — \textquotedblleft{}That depends on the time.\textquotedblright{}
\end{quote}

\subsection*{You'll want to know: \ja{なにに} \ja{よりますか}}
To ask, put \textbf{\ja{なに}} (\emph{nani}), \textquotedblleft{}what\textquotedblright{}, before \textbf{\ja{に}}:

\begin{quote}
\textbf{\ja{なにに} \ja{よりますか。}}
\end{quote}

\begin{quote}
\emph{nani ni yorimasu ka.} — \textquotedblleft{}What does it depend on?\textquotedblright{}
\end{quote}

The answer often goes on with \textbf{-\ja{たら}}:

\begin{quote}
\textbf{\ja{てんきに} \ja{よります。}}
\end{quote}

\begin{quote}
\emph{tenki ni yorimasu.} — \textquotedblleft{}It depends on the weather.\textquotedblright{}
\end{quote}

\begin{quote}
\textbf{\ja{はれたら、いきます。}}
\end{quote}

\begin{quote}
\emph{haretara, ikimasu.} — \textquotedblleft{}If it is fine, I will go.\textquotedblright{}
\end{quote}

\subsection*{Your turn}
Say these aloud:

\begin{itemize}
\raggedright
  \item \textquotedblleft{}It depends on the weather\textquotedblright{} — \emph{tenki ni yorimasu}
  \item \textquotedblleft{}It depends on the time\textquotedblright{} — \emph{jikan ni yorimasu}
  \item \textquotedblleft{}It depends on the money\textquotedblright{} — \emph{okane ni yorimasu}
\end{itemize}

\begin{tcolorbox}[breakable,colback=teal!4,colframe=teal!35!black,title={Before you move on}]
\textquotedblleft{}It depends on the weather\textquotedblright{}? (\textbf{\ja{てんきに} \ja{よります}}.) \textquotedblleft{}What does it depend on?\textquotedblright{} (\textbf{\ja{なにに} \ja{よりますか}}.) Which particle marks what it depends on? (\textbf{\ja{に}}.)
\end{tcolorbox}
\section[\texorpdfstring{\ja{ばあい} (baai)}{ばあい (baai)}]{\ja{ばあい} (baai) — a case; an occasion}

\label{lesson:JA-C261-baai}



\begin{quote}
Before the new word: ask \textquotedblleft{}what does it depend on?\textquotedblright{} (\emph{nani ni yorimasu ka}.) Answer \textquotedblleft{}it depends on the weather\textquotedblright{}. (\emph{tenki ni yorimasu}.) Then say \textquotedblleft{}calling off\textquotedblright{}. (\emph{chūshi}.)
\end{quote}

\subsection*{You'll want to know: \ja{ばあい}}
\textbf{\ja{ばあい}} — \emph{baai} — \textquotedblleft{}a case\textquotedblright{}, \textquotedblleft{}an occasion\textquotedblright{}: what is done when something happens. A noun before it takes \textbf{\ja{の}}:

\begin{quote}
\textbf{\ja{あめの} \ja{ばあいは、ちゅうしです。}}
\end{quote}

\begin{quote}
\emph{ame no baai wa, chūshi desu.} — \textquotedblleft{}In case of rain, it is called off.\textquotedblright{}
\end{quote}

\begin{quote}
\textbf{\ja{あめの} \ja{ばあいは、まつりは} \ja{ちゅうしです。}}
\end{quote}

\begin{quote}
\emph{ame no baai wa, matsuri wa chūshi desu.} — \textquotedblleft{}In case of rain, the festival is called off.\textquotedblright{}
\end{quote}

\begin{cousinweb}
\textbf{\ja{ばあい}} (\emph{baai}) is a native Japanese word (wago), inherited rather than borrowed from Chinese. JMdict writes it with two kanji that KANJIDIC2 reads the native way: \ja{ば}, \textquotedblleft{}place\textquotedblright{}, and \ja{あい}, a reading of a kanji KANJIDIC2 glosses \textquotedblleft{}fit, suit, join\textquotedblright{}. JMdict glosses the word \textquotedblleft{}case; occasion; situation\textquotedblright{} (JMdict, KANJIDIC2).
\end{cousinweb}

\subsection*{Your turn}
Say these aloud:

\begin{itemize}
\raggedright
  \item \emph{baai}
  \item \textquotedblleft{}In case of rain, it is called off\textquotedblright{} — \emph{ame no baai wa, chūshi desu}
  \item \textquotedblleft{}In case of snow, it is called off\textquotedblright{} — \emph{yuki no baai wa, chūshi desu}
\end{itemize}

\begin{tcolorbox}[breakable,colback=teal!4,colframe=teal!35!black,title={Before you move on}]
\textquotedblleft{}A case\textquotedblright{}? (\textbf{\ja{ばあい}}, \emph{baai}.) \textquotedblleft{}In case of rain\textquotedblright{}? (\textbf{\ja{あめの} \ja{ばあいは}}.)
\end{tcolorbox}
\section[\texorpdfstring{\ja{〜}\ja{たら} … \ja{でしょう} (tara … deshō)}{〜たら … でしょう (tara … deshō)}]{\ja{〜}\ja{たら} … \ja{でしょう} (tara … deshō) — what would happen}

\label{lesson:JA-C261-kau-deshou}



\begin{quote}
Say \textquotedblleft{}if I have time, I will see a film\textquotedblright{}. (\emph{jikan ga attara, eiga o mimasu}.) Then \textquotedblleft{}it will probably rain tomorrow\textquotedblright{}. (\emph{ashita ame ga furu deshō}.)
\end{quote}

\subsection*{You'll want to know: -\ja{たら} … \ja{でしょう}}
\begin{quote}
\textbf{\ja{おかねが} \ja{あったら、ちちは} \ja{かうでしょう。}}
\end{quote}

\begin{quote}
\emph{okane ga attara, chichi wa kau deshō.} — \textquotedblleft{}If he had the money, my father would buy it.\textquotedblright{}
\end{quote}

Japanese has no separate form for \textquotedblleft{}would\textquotedblright{}. The \textbf{-\ja{たら}} clause supposes, and \textbf{\ja{でしょう}} makes the result a guess: \textquotedblleft{}if there were money, he would probably buy it\textquotedblright{}. The same sentence, said about a real plan, is \textquotedblleft{}if he has the money, he will buy it\textquotedblright{}; the situation tells you which.

\begin{quote}
\textbf{\ja{じかんが} \ja{あったら、ははは} \ja{くるでしょう。}}
\end{quote}

\begin{quote}
\emph{jikan ga attara, haha wa kuru deshō.} — \textquotedblleft{}If she had time, my mother would come.\textquotedblright{}
\end{quote}

\subsection*{You'll want to know: -\ja{たら} … \ja{と} \ja{おもいます}}
For yourself, \textbf{\ja{と} \ja{おもいます}} does the same job:

\begin{quote}
\textbf{\ja{じかんが} \ja{あったら、いくと} \ja{おもいます。}}
\end{quote}

\begin{quote}
\emph{jikan ga attara, iku to omoimasu.} — \textquotedblleft{}If I had time, I think I would go.\textquotedblright{}
\end{quote}

\begin{quote}
\textbf{\ja{おかねが} \ja{あったら、かうと} \ja{おもいます。}}
\end{quote}

\begin{quote}
\emph{okane ga attara, kau to omoimasu.} — \textquotedblleft{}If I had the money, I think I would buy it.\textquotedblright{}
\end{quote}

What would have happened in the past, said with regret, needs other forms, which come later.

\subsection*{Your turn}
Say these aloud:

\begin{itemize}
\raggedright
  \item \textquotedblleft{}If I had time, I think I would go\textquotedblright{} — \emph{jikan ga attara, iku to omoimasu}
  \item \textquotedblleft{}If I had time, I think I would rest\textquotedblright{} — \emph{jikan ga attara, yasumu to omoimasu}
  \item \textquotedblleft{}If I had time, I think I would swim\textquotedblright{} — \emph{jikan ga attara, oyogu to omoimasu}
\end{itemize}

\begin{tcolorbox}[breakable,colback=teal!4,colframe=teal!35!black,title={Before you move on}]
\textquotedblleft{}He would buy it\textquotedblright{}, with the condition? (\textbf{\ja{おかねが} \ja{あったら、かうでしょう}}.) For yourself? (\textbf{\ja{かうと} \ja{おもいます}}.) Is there a separate \textquotedblleft{}would\textquotedblright{} form? (No: \textbf{-\ja{たら}} and a guess.)
\end{tcolorbox}
\section[(six lines)]{Six lines, one trip}

\label{lesson:JA-W261-six-lines}



\begin{quote}
Four lines once held together by time words and by the person you left out. Six lines can hold together by the joints you now own: \textbf{-\ja{て}}, \textbf{\ja{それから}}, \textbf{\ja{から}} and \textbf{\ja{ので}}, \textbf{-\ja{たら}}.

The day at the river was thirteen lines of this kind. Pick one \textbf{\ja{ともだち}} who will read about a trip you took.
\end{quote}

\subsection*{What to know first}
Each line has a job, and a frame to hang it on:

\begin{enumerate}
\raggedright
  \item When and where, and what was going on: \textbf{\ja{きのう}} or \textbf{\ja{きょねん}}, a place, and \textbf{-\ja{ていました}} (\textbf{\ja{あめが} \ja{ふっていました}}).
  \item The first thing that happened: \textbf{-\ja{て}}, or \textbf{\ja{それから}} at the start.
  \item A problem: \textbf{\ja{のりおくれました}}, \textbf{\ja{なくしました}}, or \textbf{\ja{こわれました}}.
  \item Why: a reason with \textbf{\ja{から}} or \textbf{\ja{ので}}.
  \item How it felt: \textbf{-\ja{かったです}}, or \textbf{\ja{でした}} after a na-adjective; \textbf{\ja{いちばん}} if one thing stood out.
  \item What you would do next time: \textbf{\ja{もし}} … \textbf{-\ja{たら}} … \textbf{\ja{と} \ja{おもいます}}.
\end{enumerate}

Use at least two of the six kanji you have copied: \textbf{\ja{月日}}, \textbf{\ja{年}\ja{じゅう}}, \textbf{\ja{言}\ja{い}\ja{出}\ja{す}}, \textbf{\ja{日本人}}, \textbf{\ja{日本中}}, \textbf{\ja{今}\ja{まで}}. Every other word goes in hiragana, spaced as the book spaces it.

\subsection*{Writing — connected composition}
Write \textbf{six lines} to the \textbf{\ja{ともだち}} you picked, one job a line, in that order. There is no model to copy: the frames are the whole of the help. Count your signs when you finish; about one hundred is a good length.

End every line with \textbf{\ja{。}}. A line may open with \textbf{\ja{それから}} or \textbf{\ja{そして}}, and then \textbf{\ja{、}} follows it.

\begin{tcolorbox}[breakable,colback=teal!4,colframe=teal!35!black,title={Before you move on}]
Read your fourth line alone. Does it say why, with \textbf{\ja{から}} or \textbf{\ja{ので}}? Does your third line say what went wrong? If each line does its one job, the six lean on each other, and that is what connected means.
\end{tcolorbox}
\section[(review)]{Review — it depends, and what would happen}

\label{lesson:JA-R261-it-depends}



\begin{quote}
Say \textquotedblleft{}it depends on the weather\textquotedblright{}. (\emph{tenki ni yorimasu}.) Then \textquotedblleft{}in case of rain, it is called off\textquotedblright{}. (\emph{ame no baai wa, chūshi desu}.)
\end{quote}

\subsection*{Your turn}
Say these aloud:

\begin{itemize}
\raggedright
  \item That depends on the time.
  \item In case of snow, the festival is called off.
  \item If I had time, I think I would go.
  \item If she had time, my mother would come.
  \item If it is fine, I will go.
  \item If it snows, what will you do?
  \item It depends on the time. I work from nine to five.
\end{itemize}

\begin{tcolorbox}[breakable,colback=teal!4,colframe=teal!35!black,title={Before you move on}]
\textquotedblleft{}It depends on\textquotedblright{}? (\textbf{\ja{〜}\ja{に} \ja{よります}}.) \textquotedblleft{}A case\textquotedblright{}? (\textbf{\ja{ばあい}}.) \textquotedblleft{}He would buy it\textquotedblright{}? (\textbf{\ja{かうでしょう}}.) \textquotedblleft{}I think I would go\textquotedblright{}? (\textbf{\ja{いくと} \ja{おもいます}}.)
\end{tcolorbox}
